\scanpage{15}
à un sommet $s\in S_M$, formé des $a\in SA_M$ et $f\in SF_M$ incidents à $s$. Choisissons un $x\in\widetilde{\deleted{\unsure{F}}S}_M\deleted{\unsure{=}}\subset M^*$ tel que $s=\dot x$, alors les $\dot t\in S_M$ tels que $\{s,t\}\in SA_M$ correspondent biunivoquement aux $y\in M$ tels que l'on ait
\[
x\wedge y=e_M\tag*{(35)}
\]
et deux tels éléments $t$\inserted{$=\dot x$} et $u$\inserted{$=\dot z$}, \unsure{correspondant|correspondent} à $y,z\in M$ tels que
\[
x\wedge y=x\wedge z=e_M,\tag*{($36'$)}
\]
sont ``incidents'', i.e. forment avec $s$ une face
\[
\{s,t,u\}\unsure{\in|=}\{\dot x,\dot y,\dot z\}\in SF_M
\]
si et seulement si on peut trouver des signes $\varepsilon,\varepsilon',\varepsilon''$ tels que l'on ait
\[
\varepsilon x+\varepsilon'y+\varepsilon''z=0.
\]
Quitte à multiplier par $\varepsilon$, \unsure{OPS|\deleted{OPS}} qu'on a
\[
x+\varepsilon'y+\varepsilon''z=0
\]
ce qui implique, en multipliant par $x\wedge$
\[
x\wedge(\varepsilon'y+\varepsilon''z)=0
\]
alors qu'on a par ($36'$)
\[
x\wedge(y-z)=0.
\]
Si on avait $\varepsilon'=\varepsilon''$, on aurait \deleted{$x\wedge\unsure{(y|y}$} \inserted{aussi}
\[
x\wedge(y+z)=0
\]
% archive sheet number: 14
